\begin{table}[tb]
\centering
\caption{ALE-based feature importance and comparison with SHAP-based rankings for the best classification model (tuned XGBoost). First-order ALE was computed on a 10,000-row explanation sample with 20 bins.}
\label{tab:ale-summaries}
\begin{tabular}{lcc}
\toprule
Metric & Value \\
\midrule
Spearman's $\rho$ (ALE vs.\ SHAP importance) & 0.76 \\
Kendall's $\tau$ (ALE vs.\ SHAP importance) & 0.58 \\
$p$-value (ALE--SHAP rank correlation) & $<0.001$ \\
\midrule
\multicolumn{2}{l}{\textbf{Top-ranked predictors (by ALE)}} \\
Rank 1 & Home possessions (HOMEPOS) \\
Rank 2 & Mathematics self-efficacy (MATHEFF) \\
Rank 3 & Grade \\
\midrule
\multicolumn{2}{l}{\textbf{Notable ranking divergences (SHAP vs.\ ALE)}} \\
ICT self-efficacy (ICTEFFIC) & Higher ALE rank than SHAP rank \\
ICT resources (ICTRES) & Lower ALE rank than SHAP rank \\
\bottomrule
\end{tabular}

\par\small The top three predictors by ALE were identical to those by permutation importance, providing cross-method validation. Divergences for ICT variables likely reflect differing treatment of feature correlations: SHAP assumes independence, inflating importance of features with broader marginal distributions, while ALE conditions on local intervals.}
\end{table}
